\chapter{Experimentación}
\label{chapter:6-Experimentación}

En este capítulo se presentan los experimentos realizados para evaluar el desempeño del sistema, cuyos criterios de evaluación se establecieron en los capítulos anteriores.

\section{Evaluación de Métodos de Transcripción}
\label{section:6-Evaluación de Métodos de Transcripción}
En la siguiente sección se escribe un análisis acerca del método de transcripción.
\subsection{Proceso de Evaluación}

La evaluación de los métodos de transcripción requiere un balance entre los 3 parametros definidos en \autoref{subsection:3-RNF6: Evaluación de Herramientas de Transcripción}. Para validar el desempeño de cada método, es necesario contar con un estándar de referencia de transcripción perfecta (\textit{ground-truth}). Con este propósito, se transcribieron manualmente siguientes clases:

\begin{itemize}
\item Computación 1, Clase 1 - Teoría 2024 \href{https://open.fing.edu.uy/courses/comp1-2024/2/}{(disponible en línea)}
\item Computación 1, Clase 2 - Teoría 2024 \href{https://open.fing.edu.uy/courses/comp1-2024/1/}{(disponible en línea)}
\item Introducción al Procesamiento de Lenguaje Natural, Clase 1 - Teoría 2016 \href{https://open.fing.edu.uy/courses/pln/1}{(disponible en línea)}
\item Introducción al Procesamiento de Lenguaje Natural, Clase 2 - Teoría 2016 \href{https://open.fing.edu.uy/courses/pln/1/}{(disponible en línea)}
\end{itemize}

El conjunto de datos así construido representa aproximadamente 5 horas y media de contenido de audio en español. Este volumen de material proporciona una base empírica suficiente para evaluar la consistencia y el rendimiento de los modelos en contextos académicos variados.

\subsection{Métricas de Evaluación}
\label{subsection:6-1-Métricas}

Para comparar el desempeño de los modelos de reconocimiento automático de voz (ASR), se emplearon cinco métricas estándar en la evaluación de sistemas de transcripción. La métrica principal de análisis es la \textit{Word Error Rate} (WER), que mide el porcentaje de palabras incorrectamente transcritas. Las métricas complementarias incluyen \textit{Match Error Rate} (MER), \textit{Word Information Lost} (WIL), \textit{Word Information Preserved} (WIP) y \textit{Character Error Rate} (CER).



\begin{table}[H]
\centering


\begin{tabular}{l l l S[table-format=1.3] S[table-format=1.3] S[table-format=1.3] S[table-format=1.3] S[table-format=1.3]}
\toprule
\textbf{Modelo} & \textbf{Conjunto de Prueba} & \textbf{Idioma} & \textbf{WER} & \textbf{MER} & \textbf{WIL} & \textbf{WIP} & \textbf{CER} \\
\midrule
% --- MODELO BASE ---
base   & Comp1-1 & ad & 0.203 & 0.200 & 0.324 & 0.676 & 0.077 \\
       &         & es & 0.198 & 0.195 & 0.318 & 0.682 & 0.074 \\
\cmidrule(r){2-8}
base   & Comp1-2 & ad & 0.199 & 0.195 & 0.314 & 0.686 & 0.077 \\
       &         & es & 0.208 & 0.203 & 0.327 & 0.673 & 0.081 \\
\midrule
% --- MODELO SMALL ---
small  & Comp1-1 & ad & 0.151 & 0.149 & 0.248 & 0.752 & 0.054 \\
       &         & es & 0.151 & 0.149 & 0.248 & 0.752 & 0.054 \\
\cmidrule(r){2-8}
small  & Comp1-2 & ad & 0.139 & 0.138 & 0.220 & 0.780 & 0.060 \\
       &         & es & 0.234 & 0.228 & 0.372 & 0.628 & 0.081 \\
\cmidrule(r){2-8}
small  & Pnl-1   & ad & 0.150 & 0.148 & 0.248 & 0.751 & 0.054 \\
       &         & es & 0.150 & 0.148 & 0.248 & 0.751 & 0.054 \\
\cmidrule(r){2-8}
small  & Pnl-2   & ad & 0.139 & 0.137 & 0.219 & 0.780 & 0.059 \\
       &         & es & 0.234 & 0.228 & 0.372 & 0.627 & 0.081 \\
\midrule
% --- MODELO MEDIUM ---
medium & Comp1-1 & ad & 0.046 & 0.046 & 0.072 & 0.928 & 0.024 \\
       &         & es & 0.046 & 0.046 & 0.072 & 0.928 & 0.024 \\
\cmidrule(r){2-8}
medium & Comp1-2 & ad & 0.087 & 0.085 & 0.120 & 0.880 & 0.052 \\
       &         & es & 0.043 & 0.042 & 0.058 & 0.942 & 0.027 \\
\cmidrule(r){2-8}
medium & Pnl-1   & ad & 0.051 & 0.045 & 0.071 & 0.928 & 0.023 \\
       &         & es & 0.051 & 0.045 & 0.071 & 0.928 & 0.023 \\
\cmidrule(r){2-8}
medium & Pnl-2   & ad & 0.086 & 0.085 & 0.120 & 0.879 & 0.052 \\
       &         & es & 0.042 & 0.042 & 0.057 & 0.942 & 0.027 \\
\bottomrule
\end{tabular}
\caption{Resultados de evaluación de modelos Whisper en tareas de transcripción de contenido académico en español. Se reportan métricas de error para autodetección de idioma (ad) y especificación explícita de español (es).}
\label{table:whisper_evaluation}
\end{table}



\subsubsection{Análisis de Resultados}

Los resultados evidencian diferencias significativas en el desempeño de los tres modelos evaluados. El modelo \textit{medium} obtiene el rendimiento más favorable, con una WER de 0.046 en la Clase 1 y 0.042 en la Clase 2 cuando se especifica explícitamente el idioma español. Esta mejora respecto a la autodetección de idioma (WER de 0.086 en Clase 2) indica que la indicación explícita del idioma optimiza la precisión del modelo.

En contraste, el modelo \textit{small} presenta tasas de error sustancialmente superiores, con WER de 0.150 en Clase 1 (consistente entre ambas estrategias de detección de idioma) y un rango de 0.139 a 0.234 en Clase 2. Notablemente, en esta última clase, la especificación de español resulta en una degradación significativa del desempeño (WER de 0.234), contrario a la tendencia observada en el modelo \textit{medium}.

El modelo \textit{base} exhibe el peor desempeño general, con WER de 0.203 en Clase 1, indicando que modelos de menor capacidad acumulan mayor proporción de errores en la transcripción.

Los resultados en las clases del curso de Procesamiento de Lenguaje Natural (Pnl-1 y Pnl-2) confirman las tendencias observadas en Computación 1. El modelo \textit{medium} alcanza una WER de 0.051 en Pnl-1 (idéntica entre ambas estrategias de detección) y de 0.042 en Pnl-2 con especificación explícita del idioma, resultados equivalentes a los obtenidos en el conjunto Comp1. El modelo \textit{small} mantiene una WER de 0.150 en Pnl-1 y oscila entre 0.139 y 0.234 en Pnl-2 según la estrategia de idioma empleada, replicando el comportamiento observado en Computación 1. La consistencia de los resultados a través de dos cursos con estilos pedagógicos distintos respalda la solidez del modelo \textit{medium} como opción de transcripción para el corpus de OpenFing.

\subsection{Análisis Cualitativo}

Más allá de las métricas cuantitativas, el análisis manual de los errores de transcripción revela patrones consistentes que permiten caracterizar mejor el comportamiento de los modelos.

\subsubsection{Errores Atribuibles al Proceso de Grabación}

Una proporción significativa de los errores identificados no corresponde a limitaciones del sistema ASR, sino a deficiencias en la calidad de la grabación original, tales como cortes de micrófono, ruido de fondo ambiental y artefactos acústicos. Estos fenómenos resultan imperceptibles sin referencia visual a la sesión de clase y constituyen limitaciones del material de entrada antes que del proceso de transcripción automática.

\subsubsection{Rendimiento en Vocabulario Especializado}

A diferencia de las dificultades observadas con palabras de uso general, el modelo exhibe desempeño robusto en la transcripción de terminología técnica. Los términos disciplinarios específicos del dominio de computación y procesamiento de lenguaje natural son capturados con alta precisión. Adicionalmente, los nombres propios son transcritos correctamente en la mayoría de los casos, sugiriendo que la arquitectura del modelo posee capacidades discriminativas efectivas para entidades nombradas.

\subsubsection{Muletillas y Palabras de Relleno}

Durante la evaluación se identificó una fuente recurrente de discrepancia: las muletillas (\textit{filler words}), esto es, palabras o sonidos utilizados de forma interjeccional durante pausas cognitivas, no son capturadas de manera consistente por el sistema ASR. Este fenómeno es coherente con hallazgos reportados en la literatura científica sobre detección de palabras de relleno en sistemas de reconocimiento de voz.

Desde la perspectiva de la aplicación propuesta en este trabajo, donde las transcripciónes alimentan un sistema de recuperación de información (RAG) diseñado para generar respuestas en un asistente conversacional, la omisión de muletillas puede considerarse beneficiosa. Estos elementos lingüísticos no contribuyen contenido semántico significativo y su eliminación reduce el ruido en los documentos indexados, potencialmente mejorando la calidad de las respuestas generadas. Una cuantificación sistemática de este fenómeno constituye una línea de mejora futura para optimizar aún más el procesamiento.

\subsection{Conclusión}
\label{subsection:6-1-Conclusión}

Los resultados obtenidos permiten establecer que el modelo Whisper \textit{medium} con especificación explícita del idioma español es el más adecuado para la transcripción del corpus de OpenFing. Alcanza una WER promedio inferior a 0.05 en todos los conjuntos de prueba evaluados, lo que equivale a menos de cinco errores por cada cien palabras. Esta precisión es suficiente para que los segmentos generados sean útiles como fuente de conocimiento en el pipeline RAG: los fragmentos recuperados contienen información semánticamente correcta, incluso cuando presentan errores menores de transcripción.

El modelo \textit{small} ofrece una alternativa de menor costo computacional pero con una WER entre dos y cuatro veces mayor, lo que puede degradar la calidad de los segmentos indexados y, por extensión, la relevancia de las respuestas generadas. El modelo \textit{base} no resulta adecuado para este dominio. En consecuencia, el sistema utiliza el modelo \textit{medium} como configuración de producción.

\subsection{Justificación de faster-whisper}
\label{subsection:6-1-faster-whisper}

Para la transcripción del corpus se utilizó la biblioteca \textit{faster-whisper} en lugar de la implementación original de Whisper provista por OpenAI. La diferencia central radica en el motor de inferencia: \textit{faster-whisper} emplea CTranslate2, un motor optimizado para la ejecución de modelos transformer que reduce el consumo de memoria y el tiempo de procesamiento sin modificar la arquitectura del modelo subyacente. En la práctica, esto se traduce en una reducción del tiempo de transcripción de entre dos y cuatro veces respecto a la implementación original, manteniendo resultados equivalentes en las métricas de error. Esta eficiencia es especialmente relevante dado que el corpus completo de OpenFing comprende 4.339 clases, un volumen que haría inviable el procesamiento con herramientas de menor rendimiento.

\section{Modelo de Vectorización}
\label{section:6-Modelo de Vectorización}

\section{Pregunta Respuesta}
\label{section:6-Pregunta Respuesta}

\subsection{Modelos}
\label{subsection:6-Modelos}

Los modelos de lenguaje evaluados en los experimentos de pregunta-respuesta son \textit{meta-llama/Meta-Llama-3.1-8B-Instruct} (Llama~3.1~8B, de Meta) como modelo principal, y \textit{Qwen/Qwen2.5-72B-Instruct} como alternativa de mayor capacidad. El modelo \textit{mistralai/Mistral-7B-Instruct-v0.3} fue evaluado en etapas previas del proyecto pero descartado en favor de Llama~3.1~8B por su menor desempeño en el dominio académico en español.

\subsection{Resultados}
\label{subsection:6-QA-Resultados}

\subsection{Conclusión}
\label{subsection:6-QA-Conclusión}



